\documentclass[a4paper,11pt, oneside,titlepage]{article}
\usepackage[pdftex]{graphicx}
\usepackage[latin2]{inputenc}
\usepackage{graphics}
\usepackage{amsmath,amssymb}
\pagestyle{plain}

\setlength{\topmargin}{-30pt}
\setlength{\headheight}{12pt}
\setlength{\headsep}{12pt}
\setlength{\oddsidemargin}{-17pt}
\setlength{\textheight}{57.0pc}
\setlength{\textwidth}{40.0pc}
 

\begin{document}
%%%%%%%%%%%%%%%%%%%%%%%%%%%%%%%%%%
%%%%%%%%%%%%%%%%%%%%%%%%%%%%%%%%%%

\pagestyle{empty} 


%\noindent Title: High-order fluctuations of temperature in hot QCD matter\\

%\noindent Authors: Jinhui Chen, Wei-jie Fu, Shi Yin, and Chunjian Zhang\\[0.3ex]

\begin{center}
{\LARGE Appeal Letter for Manuscript LD19694}\\[5ex]
\end{center}
 
\noindent Dear Editorial Board member,\\[-1ex]

\noindent We sincerely thank the Editor for their thorough review of our manuscript. We are writing to respectfully appeal the decision not to send our manuscript, titled ``\textit{High-order fluctuations of temperature in hot QCD matter}'' (Manuscript ID: LD19694 Chen), out for peer review. While we fully respect the editorial decision, we would like to highligh the novelty, significance, and broader implications of our paper, which we believe merit further consideration. We greatly appreciate your time and your consideration for publication.\\

\noindent Recent advances in heavy-ion collision experiments, where the thermal fluctuations can be isolated from other effects, e.g., the initial state geometry fluctuations [ATLAS, Phys. Rev. Lett. 133, 252301 (2024)], have opened unprecedented opportunities to probe the fundamental properties of QCD matter. For instance, one is now able to not only detect shapes of atomic nuclei [STAR, Nature 635, 67 (2024)], but also measure the speed of sound in quark-gluon plasma [CMS, Rept. Prog. Phys. 87, 077801 (2024)] directly with high-energy nuclear collisions.  \\

\noindent This raises a fundamental question: Can high-order thermodynamic observables-particularly temperature fluctuations-be isolated and measured in heavy-ion collisions? Similar with the conventional high-order net-proton number fluctuations [STAR, Phys. Rev. Lett. 126, 092301 (2021); 128, 022301 (2022); 130, 082301 (2025)], temperature fluctuations could potentially be used to probe the QCD thermodynamics and QCD phase transitions, since they are more sensitive to the thermal or critical fluctuations, i.e., the singular part of the thermodynamic potential during the phase transition. Furthermore, the experimental discovery of high-order thermodynamic fluctuations in QCD matter would represent a landmark achievement in the field.\\

\noindent To successfully isolate and discover the temperature fluctuations through event-by-event mean transverse momentum fluctuation measurements of final-state charged particles that are currently pursued at major heavy-ion facilities, a thorough theoretical understanding of their thermodynamic properties must first be established. In particular, the smoking-gun signatures of temperature fluctuations need to be identified. This work provides the first systematic theoretical investigation of these crucial aspects, addressing a significant gap in the existing literature. \\

\noindent \textbf{Scientific merit of the paper and why it reached the high standard of the journal:}\\[-1.5ex] 

\begin{itemize}

\item[$\bullet$] A new thermodynamic state function is introduced to describe the thermodynamics relevant for the mean transverse momentum fluctuations of charged particles in heavy-ion collisions. This state function has never been used in heavy-ion physics.

\item[$\bullet$] Using this thermodynamic state function, we derive analytic expressions for arbitrary-order temperature fluctuations from the basic thermodynamic relations for the first time, revealing their fundamental relationship with the entropy, heat capacity, and high-order entropy fluctuations.

\item[$\bullet$] It is found that the temperature fluctuations are suppressed remarkably as the system transitions from the hadron resonance gas (HRG) to the quark-gluon plasma (QGP) with increasing temperature or baryon chemical potential. The negative skewness of thermodynamical temperature fluctuations can be regarded as a smoking-gun signature of temperature fluctuations. In contrast, the hydrodynamic contribution to the skewness of mean transverse momentum fluctuations is found to be positive [arXiv:2004.09799].

\item[$\bullet$] These findings are general and model-independent, stemming from a fundamental thermodynamic property: the heat capacity of QCD matter increases significantly in QGP compared to that in HRG. In QGP a tiny change of the temperature would cost a huge amount of energy due to the large heat capacity, so the temperature tends not to change in comparison to the case in HRG, which implies that the temperature fluctuations would be suppressed remarkably as the matter evolves from the HRG phase to the QGP phase. This suppression naturally produces negative skewness in the fluctuation distribution, as the narrowing of high-temperature fluctuations creates an asymmetric probability density favoring lower temperatures. 

\item[$\bullet$] These predictions provide a unique signature to discover the thermodynamical temperature fluctuations in upcoming heavy-ion collision experiments, paving a novel way to study QCD thermodynamics and QCD phase diagram through experimental measurements of the mean transverse momentum fluctuations of charged particles. \\[-1.5ex] 

\end{itemize}

\noindent We have carefully revised the manuscript to enhance to enhance its novelty, impact and interest more transparent and compelling. We believe this paper makes important contributions to the field and will be of wide interest. Therefore, we respectfully request that the editorial board reconsider our manuscript in light of these key contributions.  \\

\noindent Thank you for considering this appeal. We are looking forward to your response.\\


\noindent Sincerely yours, \\[0.3ex]

\noindent Jinhui Chen, Wei-jie Fu, Shi Yin, and Chunjian Zhang

\noindent 




%%%%%%%%%%%%%%%%%%%%%%%%%%%%%%%%%%
%%%%%%%%%%%%%%%%%%%%%%%%%%%%%%%%%%
\end{document}